% Options for packages loaded elsewhere
\PassOptionsToPackage{unicode}{hyperref}
\PassOptionsToPackage{hyphens}{url}
\documentclass[
]{article}
\usepackage{xcolor}
\usepackage[margin=1in]{geometry}
\usepackage{amsmath,amssymb}
\setcounter{secnumdepth}{-\maxdimen} % remove section numbering
\usepackage{iftex}
\ifPDFTeX
  \usepackage[T1]{fontenc}
  \usepackage[utf8]{inputenc}
  \usepackage{textcomp} % provide euro and other symbols
\else % if luatex or xetex
  \usepackage{unicode-math} % this also loads fontspec
  \defaultfontfeatures{Scale=MatchLowercase}
  \defaultfontfeatures[\rmfamily]{Ligatures=TeX,Scale=1}
\fi
\usepackage{lmodern}
\ifPDFTeX\else
  % xetex/luatex font selection
\fi
% Use upquote if available, for straight quotes in verbatim environments
\IfFileExists{upquote.sty}{\usepackage{upquote}}{}
\IfFileExists{microtype.sty}{% use microtype if available
  \usepackage[]{microtype}
  \UseMicrotypeSet[protrusion]{basicmath} % disable protrusion for tt fonts
}{}
\makeatletter
\@ifundefined{KOMAClassName}{% if non-KOMA class
  \IfFileExists{parskip.sty}{%
    \usepackage{parskip}
  }{% else
    \setlength{\parindent}{0pt}
    \setlength{\parskip}{6pt plus 2pt minus 1pt}}
}{% if KOMA class
  \KOMAoptions{parskip=half}}
\makeatother
\usepackage{color}
\usepackage{fancyvrb}
\newcommand{\VerbBar}{|}
\newcommand{\VERB}{\Verb[commandchars=\\\{\}]}
\DefineVerbatimEnvironment{Highlighting}{Verbatim}{commandchars=\\\{\}}
% Add ',fontsize=\small' for more characters per line
\usepackage{framed}
\definecolor{shadecolor}{RGB}{248,248,248}
\newenvironment{Shaded}{\begin{snugshade}}{\end{snugshade}}
\newcommand{\AlertTok}[1]{\textcolor[rgb]{0.94,0.16,0.16}{#1}}
\newcommand{\AnnotationTok}[1]{\textcolor[rgb]{0.56,0.35,0.01}{\textbf{\textit{#1}}}}
\newcommand{\AttributeTok}[1]{\textcolor[rgb]{0.13,0.29,0.53}{#1}}
\newcommand{\BaseNTok}[1]{\textcolor[rgb]{0.00,0.00,0.81}{#1}}
\newcommand{\BuiltInTok}[1]{#1}
\newcommand{\CharTok}[1]{\textcolor[rgb]{0.31,0.60,0.02}{#1}}
\newcommand{\CommentTok}[1]{\textcolor[rgb]{0.56,0.35,0.01}{\textit{#1}}}
\newcommand{\CommentVarTok}[1]{\textcolor[rgb]{0.56,0.35,0.01}{\textbf{\textit{#1}}}}
\newcommand{\ConstantTok}[1]{\textcolor[rgb]{0.56,0.35,0.01}{#1}}
\newcommand{\ControlFlowTok}[1]{\textcolor[rgb]{0.13,0.29,0.53}{\textbf{#1}}}
\newcommand{\DataTypeTok}[1]{\textcolor[rgb]{0.13,0.29,0.53}{#1}}
\newcommand{\DecValTok}[1]{\textcolor[rgb]{0.00,0.00,0.81}{#1}}
\newcommand{\DocumentationTok}[1]{\textcolor[rgb]{0.56,0.35,0.01}{\textbf{\textit{#1}}}}
\newcommand{\ErrorTok}[1]{\textcolor[rgb]{0.64,0.00,0.00}{\textbf{#1}}}
\newcommand{\ExtensionTok}[1]{#1}
\newcommand{\FloatTok}[1]{\textcolor[rgb]{0.00,0.00,0.81}{#1}}
\newcommand{\FunctionTok}[1]{\textcolor[rgb]{0.13,0.29,0.53}{\textbf{#1}}}
\newcommand{\ImportTok}[1]{#1}
\newcommand{\InformationTok}[1]{\textcolor[rgb]{0.56,0.35,0.01}{\textbf{\textit{#1}}}}
\newcommand{\KeywordTok}[1]{\textcolor[rgb]{0.13,0.29,0.53}{\textbf{#1}}}
\newcommand{\NormalTok}[1]{#1}
\newcommand{\OperatorTok}[1]{\textcolor[rgb]{0.81,0.36,0.00}{\textbf{#1}}}
\newcommand{\OtherTok}[1]{\textcolor[rgb]{0.56,0.35,0.01}{#1}}
\newcommand{\PreprocessorTok}[1]{\textcolor[rgb]{0.56,0.35,0.01}{\textit{#1}}}
\newcommand{\RegionMarkerTok}[1]{#1}
\newcommand{\SpecialCharTok}[1]{\textcolor[rgb]{0.81,0.36,0.00}{\textbf{#1}}}
\newcommand{\SpecialStringTok}[1]{\textcolor[rgb]{0.31,0.60,0.02}{#1}}
\newcommand{\StringTok}[1]{\textcolor[rgb]{0.31,0.60,0.02}{#1}}
\newcommand{\VariableTok}[1]{\textcolor[rgb]{0.00,0.00,0.00}{#1}}
\newcommand{\VerbatimStringTok}[1]{\textcolor[rgb]{0.31,0.60,0.02}{#1}}
\newcommand{\WarningTok}[1]{\textcolor[rgb]{0.56,0.35,0.01}{\textbf{\textit{#1}}}}
\usepackage{longtable,booktabs,array}
\newcounter{none} % for unnumbered tables
\usepackage{calc} % for calculating minipage widths
% Correct order of tables after \paragraph or \subparagraph
\usepackage{etoolbox}
\makeatletter
\patchcmd\longtable{\par}{\if@noskipsec\mbox{}\fi\par}{}{}
\makeatother
% Allow footnotes in longtable head/foot
\IfFileExists{footnotehyper.sty}{\usepackage{footnotehyper}}{\usepackage{footnote}}
\makesavenoteenv{longtable}
\usepackage{graphicx}
\makeatletter
\newsavebox\pandoc@box
\newcommand*\pandocbounded[1]{% scales image to fit in text height/width
  \sbox\pandoc@box{#1}%
  \Gscale@div\@tempa{\textheight}{\dimexpr\ht\pandoc@box+\dp\pandoc@box\relax}%
  \Gscale@div\@tempb{\linewidth}{\wd\pandoc@box}%
  \ifdim\@tempb\p@<\@tempa\p@\let\@tempa\@tempb\fi% select the smaller of both
  \ifdim\@tempa\p@<\p@\scalebox{\@tempa}{\usebox\pandoc@box}%
  \else\usebox{\pandoc@box}%
  \fi%
}
% Set default figure placement to htbp
\def\fps@figure{htbp}
\makeatother
\setlength{\emergencystretch}{3em} % prevent overfull lines
\providecommand{\tightlist}{%
  \setlength{\itemsep}{0pt}\setlength{\parskip}{0pt}}
\usepackage{bookmark}
\IfFileExists{xurl.sty}{\usepackage{xurl}}{} % add URL line breaks if available
\urlstyle{same}
% fallback for those not using the hyperref driver hyperxmp:
\makeatletter
\@ifundefined{xmpquote}{\newcommand{\xmpquote}[1]{#1}}{}
\makeatother
\hypersetup{
  pdftitle={Лабораторная работа: Импорт и профилирование данных},
  pdfauthor={Кильдибеков Данил Андреевич 02271-ДМ},
  hidelinks,
  pdfcreator={LaTeX via pandoc}}

\title{Лабораторная работа: Импорт и профилирование данных}
\author{Кильдибеков Данил Андреевич 02271-ДМ}
\date{}

\begin{document}
\maketitle

{
\setcounter{tocdepth}{2}
\tableofcontents
}
\section{1. Описание данных и
цель}\label{ux43eux43fux438ux441ux430ux43dux438ux435-ux434ux430ux43dux43dux44bux445-ux438-ux446ux435ux43bux44c}

Цель-импортировать грязные наборы данных, привести типы, выявить
проблемы качества и составить план очистки. Исходные файлы в
\texttt{data/raw/} не редактируются.

\section{2. Изучение сырого
файла}\label{ux438ux437ux443ux447ux435ux43dux438ux435-ux441ux44bux440ux43eux433ux43e-ux444ux430ux439ux43bux430}

\begin{Shaded}
\begin{Highlighting}[]
\FunctionTok{read\_lines}\NormalTok{(}\StringTok{"data/raw/AirQualityUCI.csv"}\NormalTok{, }\AttributeTok{n\_max =} \DecValTok{5}\NormalTok{)}
\end{Highlighting}
\end{Shaded}

\begin{verbatim}
## [1] "Date;Time;CO(GT);PT08.S1(CO);NMHC(GT);C6H6(GT);PT08.S2(NMHC);NOx(GT);PT08.S3(NOx);NO2(GT);PT08.S4(NO2);PT08.S5(O3);T;RH;AH;;"
## [2] "10/03/2004;18.00.00;2,6;1360;150;11,9;1046;166;1056;113;1692;1268;13,6;48,9;0,7578;;"                                        
## [3] "10/03/2004;19.00.00;2;1292;112;9,4;955;103;1174;92;1559;972;13,3;47,7;0,7255;;"                                              
## [4] "10/03/2004;20.00.00;2,2;1402;88;9,0;939;131;1140;114;1555;1074;11,9;54,0;0,7502;;"                                           
## [5] "10/03/2004;21.00.00;2,2;1376;80;9,2;948;172;1092;122;1584;1203;11,0;60,0;0,7867;;"
\end{verbatim}

\begin{Shaded}
\begin{Highlighting}[]
\FunctionTok{read\_lines}\NormalTok{(}\StringTok{"data/raw/AirQualityUCI.csv"}\NormalTok{, }\AttributeTok{skip =} \DecValTok{9355}\NormalTok{)}
\end{Highlighting}
\end{Shaded}

\begin{verbatim}
##   [1] "04/04/2005;12.00.00;2,4;1142;-200;12,4;1063;293;603;175;1241;1092;26,9;18,3;0,6406;;"
##   [2] "04/04/2005;13.00.00;2,1;1003;-200;9,5;961;235;702;156;1041;770;28,3;13,5;0,5139;;"   
##   [3] "04/04/2005;14.00.00;2,2;1071;-200;11,9;1047;265;654;168;1129;816;28,5;13,1;0,5028;;" 
##   [4] ";;;;;;;;;;;;;;;;"                                                                    
##   [5] ";;;;;;;;;;;;;;;;"                                                                    
##   [6] ";;;;;;;;;;;;;;;;"                                                                    
##   [7] ";;;;;;;;;;;;;;;;"                                                                    
##   [8] ";;;;;;;;;;;;;;;;"                                                                    
##   [9] ";;;;;;;;;;;;;;;;"                                                                    
##  [10] ";;;;;;;;;;;;;;;;"                                                                    
##  [11] ";;;;;;;;;;;;;;;;"                                                                    
##  [12] ";;;;;;;;;;;;;;;;"                                                                    
##  [13] ";;;;;;;;;;;;;;;;"                                                                    
##  [14] ";;;;;;;;;;;;;;;;"                                                                    
##  [15] ";;;;;;;;;;;;;;;;"                                                                    
##  [16] ";;;;;;;;;;;;;;;;"                                                                    
##  [17] ";;;;;;;;;;;;;;;;"                                                                    
##  [18] ";;;;;;;;;;;;;;;;"                                                                    
##  [19] ";;;;;;;;;;;;;;;;"                                                                    
##  [20] ";;;;;;;;;;;;;;;;"                                                                    
##  [21] ";;;;;;;;;;;;;;;;"                                                                    
##  [22] ";;;;;;;;;;;;;;;;"                                                                    
##  [23] ";;;;;;;;;;;;;;;;"                                                                    
##  [24] ";;;;;;;;;;;;;;;;"                                                                    
##  [25] ";;;;;;;;;;;;;;;;"                                                                    
##  [26] ";;;;;;;;;;;;;;;;"                                                                    
##  [27] ";;;;;;;;;;;;;;;;"                                                                    
##  [28] ";;;;;;;;;;;;;;;;"                                                                    
##  [29] ";;;;;;;;;;;;;;;;"                                                                    
##  [30] ";;;;;;;;;;;;;;;;"                                                                    
##  [31] ";;;;;;;;;;;;;;;;"                                                                    
##  [32] ";;;;;;;;;;;;;;;;"                                                                    
##  [33] ";;;;;;;;;;;;;;;;"                                                                    
##  [34] ";;;;;;;;;;;;;;;;"                                                                    
##  [35] ";;;;;;;;;;;;;;;;"                                                                    
##  [36] ";;;;;;;;;;;;;;;;"                                                                    
##  [37] ";;;;;;;;;;;;;;;;"                                                                    
##  [38] ";;;;;;;;;;;;;;;;"                                                                    
##  [39] ";;;;;;;;;;;;;;;;"                                                                    
##  [40] ";;;;;;;;;;;;;;;;"                                                                    
##  [41] ";;;;;;;;;;;;;;;;"                                                                    
##  [42] ";;;;;;;;;;;;;;;;"                                                                    
##  [43] ";;;;;;;;;;;;;;;;"                                                                    
##  [44] ";;;;;;;;;;;;;;;;"                                                                    
##  [45] ";;;;;;;;;;;;;;;;"                                                                    
##  [46] ";;;;;;;;;;;;;;;;"                                                                    
##  [47] ";;;;;;;;;;;;;;;;"                                                                    
##  [48] ";;;;;;;;;;;;;;;;"                                                                    
##  [49] ";;;;;;;;;;;;;;;;"                                                                    
##  [50] ";;;;;;;;;;;;;;;;"                                                                    
##  [51] ";;;;;;;;;;;;;;;;"                                                                    
##  [52] ";;;;;;;;;;;;;;;;"                                                                    
##  [53] ";;;;;;;;;;;;;;;;"                                                                    
##  [54] ";;;;;;;;;;;;;;;;"                                                                    
##  [55] ";;;;;;;;;;;;;;;;"                                                                    
##  [56] ";;;;;;;;;;;;;;;;"                                                                    
##  [57] ";;;;;;;;;;;;;;;;"                                                                    
##  [58] ";;;;;;;;;;;;;;;;"                                                                    
##  [59] ";;;;;;;;;;;;;;;;"                                                                    
##  [60] ";;;;;;;;;;;;;;;;"                                                                    
##  [61] ";;;;;;;;;;;;;;;;"                                                                    
##  [62] ";;;;;;;;;;;;;;;;"                                                                    
##  [63] ";;;;;;;;;;;;;;;;"                                                                    
##  [64] ";;;;;;;;;;;;;;;;"                                                                    
##  [65] ";;;;;;;;;;;;;;;;"                                                                    
##  [66] ";;;;;;;;;;;;;;;;"                                                                    
##  [67] ";;;;;;;;;;;;;;;;"                                                                    
##  [68] ";;;;;;;;;;;;;;;;"                                                                    
##  [69] ";;;;;;;;;;;;;;;;"                                                                    
##  [70] ";;;;;;;;;;;;;;;;"                                                                    
##  [71] ";;;;;;;;;;;;;;;;"                                                                    
##  [72] ";;;;;;;;;;;;;;;;"                                                                    
##  [73] ";;;;;;;;;;;;;;;;"                                                                    
##  [74] ";;;;;;;;;;;;;;;;"                                                                    
##  [75] ";;;;;;;;;;;;;;;;"                                                                    
##  [76] ";;;;;;;;;;;;;;;;"                                                                    
##  [77] ";;;;;;;;;;;;;;;;"                                                                    
##  [78] ";;;;;;;;;;;;;;;;"                                                                    
##  [79] ";;;;;;;;;;;;;;;;"                                                                    
##  [80] ";;;;;;;;;;;;;;;;"                                                                    
##  [81] ";;;;;;;;;;;;;;;;"                                                                    
##  [82] ";;;;;;;;;;;;;;;;"                                                                    
##  [83] ";;;;;;;;;;;;;;;;"                                                                    
##  [84] ";;;;;;;;;;;;;;;;"                                                                    
##  [85] ";;;;;;;;;;;;;;;;"                                                                    
##  [86] ";;;;;;;;;;;;;;;;"                                                                    
##  [87] ";;;;;;;;;;;;;;;;"                                                                    
##  [88] ";;;;;;;;;;;;;;;;"                                                                    
##  [89] ";;;;;;;;;;;;;;;;"                                                                    
##  [90] ";;;;;;;;;;;;;;;;"                                                                    
##  [91] ";;;;;;;;;;;;;;;;"                                                                    
##  [92] ";;;;;;;;;;;;;;;;"                                                                    
##  [93] ";;;;;;;;;;;;;;;;"                                                                    
##  [94] ";;;;;;;;;;;;;;;;"                                                                    
##  [95] ";;;;;;;;;;;;;;;;"                                                                    
##  [96] ";;;;;;;;;;;;;;;;"                                                                    
##  [97] ";;;;;;;;;;;;;;;;"                                                                    
##  [98] ";;;;;;;;;;;;;;;;"                                                                    
##  [99] ";;;;;;;;;;;;;;;;"                                                                    
## [100] ";;;;;;;;;;;;;;;;"                                                                    
## [101] ";;;;;;;;;;;;;;;;"                                                                    
## [102] ";;;;;;;;;;;;;;;;"                                                                    
## [103] ";;;;;;;;;;;;;;;;"                                                                    
## [104] ";;;;;;;;;;;;;;;;"                                                                    
## [105] ";;;;;;;;;;;;;;;;"                                                                    
## [106] ";;;;;;;;;;;;;;;;"                                                                    
## [107] ";;;;;;;;;;;;;;;;"                                                                    
## [108] ";;;;;;;;;;;;;;;;"                                                                    
## [109] ";;;;;;;;;;;;;;;;"                                                                    
## [110] ";;;;;;;;;;;;;;;;"                                                                    
## [111] ";;;;;;;;;;;;;;;;"                                                                    
## [112] ";;;;;;;;;;;;;;;;"                                                                    
## [113] ";;;;;;;;;;;;;;;;"                                                                    
## [114] ";;;;;;;;;;;;;;;;"                                                                    
## [115] ";;;;;;;;;;;;;;;;"                                                                    
## [116] ";;;;;;;;;;;;;;;;"                                                                    
## [117] ";;;;;;;;;;;;;;;;"
\end{verbatim}

\begin{Shaded}
\begin{Highlighting}[]
\FunctionTok{guess\_encoding}\NormalTok{(}\StringTok{"data/raw/AirQualityUCI.csv"}\NormalTok{)}
\end{Highlighting}
\end{Shaded}

\begin{verbatim}
## # A tibble: 1 x 2
##   encoding confidence
##   <chr>         <dbl>
## 1 ASCII             1
\end{verbatim}

\textbf{Параметры файла}

{\def\LTcaptype{none} % do not increment counter
\begin{longtable}[]{@{}
  >{\raggedright\arraybackslash}p{(\linewidth - 4\tabcolsep) * \real{0.3333}}
  >{\raggedright\arraybackslash}p{(\linewidth - 4\tabcolsep) * \real{0.3333}}
  >{\raggedright\arraybackslash}p{(\linewidth - 4\tabcolsep) * \real{0.3333}}@{}}
\toprule\noalign{}
\begin{minipage}[b]{\linewidth}\raggedright
Параметр
\end{minipage} & \begin{minipage}[b]{\linewidth}\raggedright
Найденное значение
\end{minipage} & \begin{minipage}[b]{\linewidth}\raggedright
Чем подтверждено
\end{minipage} \\
\midrule\noalign{}
\endhead
\bottomrule\noalign{}
\endlastfoot
Разделитель полей & \texttt{;} &
\texttt{read\_lines(...,\ n\_max\ =\ 5)} \\
Десятичный знак & \texttt{,} & значения вида \texttt{2,6} \\
Кодировка & ASCII & \texttt{guess\_encoding()} \\
Служебные строки перед таблицей & нет, 1 строка заголовка &
\texttt{read\_lines()} \\
Код пропуска & \texttt{-200} & значения \texttt{-200} в числовых
столбцах \\
Пустые строки или столбцы в конце & есть &
\texttt{read\_lines(...,\ skip\ =\ 9355)} \\
\end{longtable}
}

\section{3. Импорт основного
набора}\label{ux438ux43cux43fux43eux440ux442-ux43eux441ux43dux43eux432ux43dux43eux433ux43e-ux43dux430ux431ux43eux440ux430}

\begin{Shaded}
\begin{Highlighting}[]
\NormalTok{df\_raw }\OtherTok{\textless{}{-}} \FunctionTok{read\_csv2}\NormalTok{(}
  \StringTok{"data/raw/AirQualityUCI.csv"}\NormalTok{,}
  \AttributeTok{na =} \FunctionTok{c}\NormalTok{(}\StringTok{""}\NormalTok{, }\StringTok{"NA"}\NormalTok{, }\StringTok{"{-}200"}\NormalTok{),}
  \AttributeTok{col\_types =} \FunctionTok{cols}\NormalTok{(}
    \AttributeTok{Date =} \FunctionTok{col\_character}\NormalTok{(),}
    \AttributeTok{Time =} \FunctionTok{col\_character}\NormalTok{(),}
    \AttributeTok{.default =} \FunctionTok{col\_double}\NormalTok{()}
\NormalTok{  ),}
  \AttributeTok{locale =} \FunctionTok{locale}\NormalTok{(}\AttributeTok{decimal\_mark =} \StringTok{","}\NormalTok{, }\AttributeTok{encoding =} \StringTok{"UTF{-}8"}\NormalTok{),}
  \AttributeTok{trim\_ws =} \ConstantTok{TRUE}
\NormalTok{)}

\FunctionTok{glimpse}\NormalTok{(df\_raw)}
\end{Highlighting}
\end{Shaded}

\begin{verbatim}
## Rows: 9,471
## Columns: 17
## $ Date            <chr> "10/03/2004", "10/03/2004", "10/03/2004", "10/03/2004"~
## $ Time            <chr> "18.00.00", "19.00.00", "20.00.00", "21.00.00", "22.00~
## $ `CO(GT)`        <dbl> 2.6, 2.0, 2.2, 2.2, 1.6, 1.2, 1.2, 1.0, 0.9, 0.6, NA, ~
## $ `PT08.S1(CO)`   <dbl> 1360, 1292, 1402, 1376, 1272, 1197, 1185, 1136, 1094, ~
## $ `NMHC(GT)`      <dbl> 150, 112, 88, 80, 51, 38, 31, 31, 24, 19, 14, 8, 16, 2~
## $ `C6H6(GT)`      <dbl> 11.9, 9.4, 9.0, 9.2, 6.5, 4.7, 3.6, 3.3, 2.3, 1.7, 1.3~
## $ `PT08.S2(NMHC)` <dbl> 1046, 955, 939, 948, 836, 750, 690, 672, 609, 561, 527~
## $ `NOx(GT)`       <dbl> 166, 103, 131, 172, 131, 89, 62, 62, 45, NA, 21, 16, 3~
## $ `PT08.S3(NOx)`  <dbl> 1056, 1174, 1140, 1092, 1205, 1337, 1462, 1453, 1579, ~
## $ `NO2(GT)`       <dbl> 113, 92, 114, 122, 116, 96, 77, 76, 60, NA, 34, 28, 48~
## $ `PT08.S4(NO2)`  <dbl> 1692, 1559, 1555, 1584, 1490, 1393, 1333, 1333, 1276, ~
## $ `PT08.S5(O3)`   <dbl> 1268, 972, 1074, 1203, 1110, 949, 733, 730, 620, 501, ~
## $ T               <dbl> 13.6, 13.3, 11.9, 11.0, 11.2, 11.2, 11.3, 10.7, 10.7, ~
## $ RH              <dbl> 48.9, 47.7, 54.0, 60.0, 59.6, 59.2, 56.8, 60.0, 59.7, ~
## $ AH              <dbl> 0.7578, 0.7255, 0.7502, 0.7867, 0.7888, 0.7848, 0.7603~
## $ ...16           <dbl> NA, NA, NA, NA, NA, NA, NA, NA, NA, NA, NA, NA, NA, NA~
## $ ...17           <dbl> NA, NA, NA, NA, NA, NA, NA, NA, NA, NA, NA, NA, NA, NA~
\end{verbatim}

\begin{Shaded}
\begin{Highlighting}[]
\FunctionTok{spec}\NormalTok{(df\_raw)}
\end{Highlighting}
\end{Shaded}

\begin{verbatim}
## cols(
##   .default = col_double(),
##   Date = col_character(),
##   Time = col_character(),
##   `CO(GT)` = col_double(),
##   `PT08.S1(CO)` = col_double(),
##   `NMHC(GT)` = col_double(),
##   `C6H6(GT)` = col_double(),
##   `PT08.S2(NMHC)` = col_double(),
##   `NOx(GT)` = col_double(),
##   `PT08.S3(NOx)` = col_double(),
##   `NO2(GT)` = col_double(),
##   `PT08.S4(NO2)` = col_double(),
##   `PT08.S5(O3)` = col_double(),
##   T = col_double(),
##   RH = col_double(),
##   AH = col_double(),
##   ...16 = col_double(),
##   ...17 = col_double()
## )
\end{verbatim}

\begin{Shaded}
\begin{Highlighting}[]
\FunctionTok{problems}\NormalTok{(df\_raw)}
\end{Highlighting}
\end{Shaded}

\begin{verbatim}
## # A tibble: 0 x 5
## # i 5 variables: row <int>, col <int>, expected <chr>, actual <chr>, file <chr>
\end{verbatim}

Пояснения: - \texttt{read\_csv2()} --- разделитель \texttt{;},
десятичная запятая - \texttt{na\ =\ c("",\ "NA",\ "-200")} --- иначе
\texttt{-200} попадёт в статистику как реальное значение - \texttt{Date}
и \texttt{Time} читаются как текст, чтобы потом собрать \texttt{POSIXct}
- \texttt{.default\ =\ col\_double()} --- числовые датчики сразу
получают числовой тип - \texttt{problems()} показывает проблемы парсинга

\section{4. Файл с проблемной
кодировкой}\label{ux444ux430ux439ux43b-ux441-ux43fux440ux43eux431ux43bux435ux43cux43dux43eux439-ux43aux43eux434ux438ux440ux43eux432ux43aux43eux439}

\subsection{4.1. Чтение по умолчанию ---
«до»}\label{ux447ux442ux435ux43dux438ux435-ux43fux43e-ux443ux43cux43eux43bux447ux430ux43dux438ux44e-ux434ux43e}

\begin{Shaded}
\begin{Highlighting}[]
\NormalTok{variant\_file }\OtherTok{\textless{}{-}} \StringTok{"data/raw/students.csv"}

\NormalTok{df\_bad }\OtherTok{\textless{}{-}} \FunctionTok{read\_delim}\NormalTok{(variant\_file, }\AttributeTok{delim =} \StringTok{";"}\NormalTok{, }\AttributeTok{show\_col\_types =} \ConstantTok{FALSE}\NormalTok{)}
\FunctionTok{head}\NormalTok{(df\_bad, }\DecValTok{5}\NormalTok{)}
\end{Highlighting}
\end{Shaded}

\begin{verbatim}
## # A tibble: 5 x 5
##   `ФÐ\u0098Ð\u009e`               Ð\u0093Ñ\u0080Ñ\u008~1 Ð\u0094иÑ\u0081Ñ\u0~2
##   <chr>                            <chr>                  <chr>                 
## 1 "СидоÑ\u0080ов Ð\u0090нÐ~ "Ð\u0098Ð\u0092Т-21"  "Ð\u0090нализ д~
## 2 "Ð\u009dиколаева Ð\u009~ "Ð\u0098Ð\u0092Т-21"  "Ð\u009cаÑ\u0082ем~
## 3 "Ð\u009aÑ\u0083знеÑ\u0086оÐ~ "Ð\u0098Ð\u0092Т-21"  "Ð\u009cаÑ\u0088ин~
## 4 "Ð\u009aÑ\u0083знеÑ\u0086оÐ~ "Ð\u0098Ð\u0092Т-22"  "Ð\u009cаÑ\u0088ин~
## 5 "Ð\u0098ванов Ð\u009cаÑ\u~ "Ð\u0098Ð\u0092Т-22"  "Ð\u0090нализ д~
## # i abbreviated names: 1: `Ð\u0093Ñ\u0080Ñ\u0083ппа`,
## #   2: `Ð\u0094иÑ\u0081Ñ\u0086иплина`
## # i 2 more variables: `Ð\u0091аллÑ\u008b` <dbl>,
## #   `Ð\u009eÑ\u0086енка` <chr>
\end{verbatim}

\subsection{4.2.
Диагностика}\label{ux434ux438ux430ux433ux43dux43eux441ux442ux438ux43aux430}

\begin{Shaded}
\begin{Highlighting}[]
\FunctionTok{readBin}\NormalTok{(variant\_file, }\StringTok{"raw"}\NormalTok{, }\DecValTok{12}\NormalTok{)}
\end{Highlighting}
\end{Shaded}

\begin{verbatim}
##  [1] c3 90 c2 a4 c3 90 c2 98 c3 90 c2 9e
\end{verbatim}

\begin{Shaded}
\begin{Highlighting}[]
\FunctionTok{guess\_encoding}\NormalTok{(variant\_file)}
\end{Highlighting}
\end{Shaded}

\begin{verbatim}
## # A tibble: 1 x 2
##   encoding confidence
##   <chr>         <dbl>
## 1 UTF-8             1
\end{verbatim}

\begin{itemize}
\tightlist
\item
  \texttt{readBin} →
  \texttt{c3\ 90\ c2\ a4\ c3\ 90\ c2\ 98\ c3\ 90\ c2\ 9e\ …}
\item
  \texttt{guess\_encoding} - UTF-8, confidence = 1
\end{itemize}

Три байта \texttt{c3\ 90\ c2\ a4} вместо одного \texttt{d0} --- признак
\textbf{двойного кодирования}: исходный UTF-8 прочитан как Latin-1 и
сохранён снова как UTF-8. Поэтому \texttt{encoding\ =\ "UTF-8"} не
помогае

\subsection{\texorpdfstring{4.3. Чтение с подобранными \texttt{skip} и
\texttt{comment} +
восстановление}{4.3. Чтение с подобранными skip и comment + восстановление}}\label{ux447ux442ux435ux43dux438ux435-ux441-ux43fux43eux434ux43eux431ux440ux430ux43dux43dux44bux43cux438-skip-ux438-comment-ux432ux43eux441ux441ux442ux430ux43dux43eux432ux43bux435ux43dux438ux435}

\begin{Shaded}
\begin{Highlighting}[]
\NormalTok{df\_fixed }\OtherTok{\textless{}{-}} \FunctionTok{read\_delim}\NormalTok{(}
\NormalTok{  variant\_file,}
  \AttributeTok{delim          =} \StringTok{";"}\NormalTok{,}
  \AttributeTok{skip           =} \DecValTok{0}\NormalTok{,}
  \AttributeTok{comment        =} \StringTok{""}\NormalTok{,}
  \AttributeTok{locale         =} \FunctionTok{locale}\NormalTok{(}\AttributeTok{encoding =} \StringTok{"UTF{-}8"}\NormalTok{),}
  \AttributeTok{show\_col\_types =} \ConstantTok{TRUE}
\NormalTok{) }\SpecialCharTok{\%\textgreater{}\%}
  \FunctionTok{rename\_with}\NormalTok{(fix\_mojibake) }\SpecialCharTok{\%\textgreater{}\%}
  \FunctionTok{mutate}\NormalTok{(}\FunctionTok{across}\NormalTok{(}\FunctionTok{where}\NormalTok{(is.character), fix\_mojibake))}

\FunctionTok{names}\NormalTok{(df\_fixed)}
\end{Highlighting}
\end{Shaded}

\begin{verbatim}
## [1] "ФИО"        "Группа"     "Дисциплина" "Баллы"      "Оценка"
\end{verbatim}

\begin{Shaded}
\begin{Highlighting}[]
\FunctionTok{head}\NormalTok{(df\_fixed, }\DecValTok{5}\NormalTok{)}
\end{Highlighting}
\end{Shaded}

\begin{verbatim}
## # A tibble: 5 x 5
##   ФИО                            Группа Дисциплина                Баллы Оценка  
##   <chr>                          <chr>  <chr>                     <dbl> <chr>   
## 1 Сидоров Андрей Сергеевич       ИВТ-21 Анализ данных                69 удовлет~
## 2 Николаева Дмитрий Сергеевич    ИВТ-21 Математическая статистика    86 отлично 
## 3 Кузнецова Екатерина Алексеевна ИВТ-21 Машинное обучение            85 отлично 
## 4 Кузнецова Алексей Алексеевна   ИВТ-22 Машинное обучение            73 хорошо  
## 5 Иванов Мария Петрович          ИВТ-22 Анализ данных                89 отлично
\end{verbatim}

\begin{Shaded}
\begin{Highlighting}[]
\FunctionTok{glimpse}\NormalTok{(df\_fixed)}
\end{Highlighting}
\end{Shaded}

\begin{verbatim}
## Rows: 36
## Columns: 5
## $ ФИО        <chr> "Сидоров Андрей Сергеевич", "Николаева Дмитрий Сергеевич", ~
## $ Группа     <chr> "ИВТ-21", "ИВТ-21", "ИВТ-21", "ИВТ-22", "ИВТ-22", "ИВТ-22",~
## $ Дисциплина <chr> "Анализ данных", "Математическая статистика", "Машинное обу~
## $ Баллы      <dbl> 69, 86, 85, 73, 89, 55, 50, 64, 80, 56, 68, 58, 63, 51, 93,~
## $ Оценка     <chr> "удовлетворительно", "отлично", "отлично", "хорошо", "отлич~
\end{verbatim}

\begin{Shaded}
\begin{Highlighting}[]
\FunctionTok{dim}\NormalTok{(df\_fixed)}
\end{Highlighting}
\end{Shaded}

\begin{verbatim}
## [1] 36  5
\end{verbatim}

\texttt{fix\_mojibake()} (в \texttt{setup}):
\texttt{iconv(x,\ from\ =\ "UTF-8",\ to\ =\ "latin1")} +
\texttt{Encoding(y)\ -\ "UTF-8"} --- обратная операция к той, что
преобразовала файл

\subsection{4.4. Подобранные
параметры}\label{ux43fux43eux434ux43eux431ux440ux430ux43dux43dux44bux435-ux43fux430ux440ux430ux43cux435ux442ux440ux44b}

{\def\LTcaptype{none} % do not increment counter
\begin{longtable}[]{@{}
  >{\raggedright\arraybackslash}p{(\linewidth - 4\tabcolsep) * \real{0.3333}}
  >{\raggedright\arraybackslash}p{(\linewidth - 4\tabcolsep) * \real{0.3333}}
  >{\raggedright\arraybackslash}p{(\linewidth - 4\tabcolsep) * \real{0.3333}}@{}}
\toprule\noalign{}
\begin{minipage}[b]{\linewidth}\raggedright
Параметр
\end{minipage} & \begin{minipage}[b]{\linewidth}\raggedright
Значение
\end{minipage} & \begin{minipage}[b]{\linewidth}\raggedright
Чем подтверждено
\end{minipage} \\
\midrule\noalign{}
\endhead
\bottomrule\noalign{}
\endlastfoot
\texttt{delim} & \texttt{";"} &
\texttt{read\_lines(...,\ n\_max\ =\ 1)} \\
\texttt{skip} & \texttt{0} & строка 1 --- заголовок \\
\texttt{comment} & \texttt{""} & нет \texttt{\#}-строк \\
\texttt{encoding} & \texttt{"UTF-8"} & \texttt{guess\_encoding()} =
UTF-8, 1; \texttt{readBin} --- валидные UTF-8-байты \\
Размер & 36 × 5 & \texttt{dim(df\_fixed)} \\
Типы & chr, chr, chr, dbl, chr & \texttt{glimpse(df\_fixed)} \\
\end{longtable}
}

\textbf{Вывод.} \texttt{delim}, \texttt{skip}, \texttt{comment}
подобраны штатно. Кодировка файла --- UTF-8, но содержимое --- двойной
mojibake, поэтому \texttt{encoding} сам по себе кириллицу не
восстанавливает. Кириллица восстановлена обратной операцией
\texttt{fix\_mojibake()}: 36 × 5, типы
\texttt{chr,\ chr,\ chr,\ dbl,\ chr}, \texttt{nchar(ФИО{[}1{]})\ =\ 24}.

\section{5. Имена, пустые строки и
типы}\label{ux438ux43cux435ux43dux430-ux43fux443ux441ux442ux44bux435-ux441ux442ux440ux43eux43aux438-ux438-ux442ux438ux43fux44b}

\begin{Shaded}
\begin{Highlighting}[]
\NormalTok{df }\OtherTok{\textless{}{-}}\NormalTok{ df\_raw }\SpecialCharTok{\%\textgreater{}\%}
  \FunctionTok{clean\_names}\NormalTok{() }\SpecialCharTok{\%\textgreater{}\%}
  \FunctionTok{remove\_empty}\NormalTok{(}\AttributeTok{which =} \FunctionTok{c}\NormalTok{(}\StringTok{"rows"}\NormalTok{, }\StringTok{"cols"}\NormalTok{)) }\SpecialCharTok{\%\textgreater{}\%}
  \FunctionTok{mutate}\NormalTok{(}
    \AttributeTok{time =} \FunctionTok{str\_replace\_all}\NormalTok{(time, }\StringTok{"}\SpecialCharTok{\textbackslash{}\textbackslash{}}\StringTok{."}\NormalTok{, }\StringTok{":"}\NormalTok{),}
    \AttributeTok{datetime =} \FunctionTok{parse\_date\_time}\NormalTok{(}
      \FunctionTok{paste}\NormalTok{(date, time),}
      \AttributeTok{orders =} \FunctionTok{c}\NormalTok{(}\StringTok{"dmy HMS"}\NormalTok{, }\StringTok{"dmy HM"}\NormalTok{),}
      \AttributeTok{tz =} \StringTok{"Europe/Rome"}
\NormalTok{    )}
\NormalTok{  )}

\FunctionTok{glimpse}\NormalTok{(df)}
\end{Highlighting}
\end{Shaded}

\begin{verbatim}
## Rows: 9,357
## Columns: 16
## $ date         <chr> "10/03/2004", "10/03/2004", "10/03/2004", "10/03/2004", "~
## $ time         <chr> "18:00:00", "19:00:00", "20:00:00", "21:00:00", "22:00:00~
## $ co_gt        <dbl> 2.6, 2.0, 2.2, 2.2, 1.6, 1.2, 1.2, 1.0, 0.9, 0.6, NA, 0.7~
## $ pt08_s1_co   <dbl> 1360, 1292, 1402, 1376, 1272, 1197, 1185, 1136, 1094, 101~
## $ nmhc_gt      <dbl> 150, 112, 88, 80, 51, 38, 31, 31, 24, 19, 14, 8, 16, 29, ~
## $ c6h6_gt      <dbl> 11.9, 9.4, 9.0, 9.2, 6.5, 4.7, 3.6, 3.3, 2.3, 1.7, 1.3, 1~
## $ pt08_s2_nmhc <dbl> 1046, 955, 939, 948, 836, 750, 690, 672, 609, 561, 527, 5~
## $ n_ox_gt      <dbl> 166, 103, 131, 172, 131, 89, 62, 62, 45, NA, 21, 16, 34, ~
## $ pt08_s3_n_ox <dbl> 1056, 1174, 1140, 1092, 1205, 1337, 1462, 1453, 1579, 170~
## $ no2_gt       <dbl> 113, 92, 114, 122, 116, 96, 77, 76, 60, NA, 34, 28, 48, 8~
## $ pt08_s4_no2  <dbl> 1692, 1559, 1555, 1584, 1490, 1393, 1333, 1333, 1276, 123~
## $ pt08_s5_o3   <dbl> 1268, 972, 1074, 1203, 1110, 949, 733, 730, 620, 501, 445~
## $ t            <dbl> 13.6, 13.3, 11.9, 11.0, 11.2, 11.2, 11.3, 10.7, 10.7, 10.~
## $ rh           <dbl> 48.9, 47.7, 54.0, 60.0, 59.6, 59.2, 56.8, 60.0, 59.7, 60.~
## $ ah           <dbl> 0.7578, 0.7255, 0.7502, 0.7867, 0.7888, 0.7848, 0.7603, 0~
## $ datetime     <dttm> 2004-03-10 18:00:00, 2004-03-10 19:00:00, 2004-03-10 20:~
\end{verbatim}

\begin{Shaded}
\begin{Highlighting}[]
\FunctionTok{sum}\NormalTok{(}\FunctionTok{is.na}\NormalTok{(df}\SpecialCharTok{$}\NormalTok{datetime))}
\end{Highlighting}
\end{Shaded}

\begin{verbatim}
## [1] 2
\end{verbatim}

\begin{Shaded}
\begin{Highlighting}[]
\NormalTok{df }\SpecialCharTok{\%\textgreater{}\%}
  \FunctionTok{filter}\NormalTok{(}\FunctionTok{is.na}\NormalTok{(datetime)) }\SpecialCharTok{\%\textgreater{}\%}
  \FunctionTok{count}\NormalTok{(date, time, }\AttributeTok{sort =} \ConstantTok{TRUE}\NormalTok{) }\SpecialCharTok{\%\textgreater{}\%}
  \FunctionTok{head}\NormalTok{(}\DecValTok{10}\NormalTok{)}
\end{Highlighting}
\end{Shaded}

\begin{verbatim}
## # A tibble: 2 x 3
##   date       time         n
##   <chr>      <chr>    <int>
## 1 27/03/2005 02:00:00     1
## 2 28/03/2004 02:00:00     1
\end{verbatim}

\textbf{Имя до --- имя после}

{\def\LTcaptype{none} % do not increment counter
\begin{longtable}[]{@{}ll@{}}
\toprule\noalign{}
До & После \\
\midrule\noalign{}
\endhead
\bottomrule\noalign{}
\endlastfoot
\texttt{CO(GT)} & \texttt{co\_gt} \\
\texttt{PT08.S1(CO)} & \texttt{pt08\_s1\_co} \\
\texttt{NMHC(GT)} & \texttt{nmhc\_gt} \\
\texttt{C6H6(GT)} & \texttt{c6h6\_gt} \\
\texttt{PT08.S2(NMHC)} & \texttt{pt08\_s2\_nmhc} \\
\texttt{NOx(GT)} & \texttt{nox\_gt} \\
\texttt{PT08.S3(NOx)} & \texttt{pt08\_s3\_nox} \\
\texttt{NO2(GT)} & \texttt{no2\_gt} \\
\texttt{PT08.S4(NO2)} & \texttt{pt08\_s4\_no2} \\
\texttt{PT08.S5(O3)} & \texttt{pt08\_s5\_o3} \\
\texttt{T} & \texttt{t} \\
\texttt{RH} & \texttt{rh} \\
\texttt{AH} & \texttt{ah} \\
\end{longtable}
}

NA в \texttt{datetime} объясняются пустыми/некорректными датой или
временем; после \texttt{remove\_empty()} и замены \texttt{.} на
\texttt{:} их должно быть немного или 0.

\section{6.
Профилирование}\label{ux43fux440ux43eux444ux438ux43bux438ux440ux43eux432ux430ux43dux438ux435}

\begin{Shaded}
\begin{Highlighting}[]
\FunctionTok{skim}\NormalTok{(df)}
\end{Highlighting}
\end{Shaded}

\begin{longtable}[]{@{}ll@{}}
\caption{Data summary}\tabularnewline
\toprule\noalign{}
\endfirsthead
\endhead
\bottomrule\noalign{}
\endlastfoot
Name & df \\
Number of rows & 9357 \\
Number of columns & 16 \\
\_\_\_\_\_\_\_\_\_\_\_\_\_\_\_\_\_\_\_\_\_\_\_ & \\
Column type frequency: & \\
character & 2 \\
numeric & 13 \\
POSIXct & 1 \\
\_\_\_\_\_\_\_\_\_\_\_\_\_\_\_\_\_\_\_\_\_\_\_\_ & \\
Group variables & None \\
\end{longtable}

\textbf{Variable type: character}

{\def\LTcaptype{none} % do not increment counter
\begin{longtable}[]{@{}
  >{\raggedright\arraybackslash}p{(\linewidth - 14\tabcolsep) * \real{0.1944}}
  >{\raggedleft\arraybackslash}p{(\linewidth - 14\tabcolsep) * \real{0.1389}}
  >{\raggedleft\arraybackslash}p{(\linewidth - 14\tabcolsep) * \real{0.1944}}
  >{\raggedleft\arraybackslash}p{(\linewidth - 14\tabcolsep) * \real{0.0556}}
  >{\raggedleft\arraybackslash}p{(\linewidth - 14\tabcolsep) * \real{0.0556}}
  >{\raggedleft\arraybackslash}p{(\linewidth - 14\tabcolsep) * \real{0.0833}}
  >{\raggedleft\arraybackslash}p{(\linewidth - 14\tabcolsep) * \real{0.1250}}
  >{\raggedleft\arraybackslash}p{(\linewidth - 14\tabcolsep) * \real{0.1528}}@{}}
\toprule\noalign{}
\begin{minipage}[b]{\linewidth}\raggedright
skim\_variable
\end{minipage} & \begin{minipage}[b]{\linewidth}\raggedleft
n\_missing
\end{minipage} & \begin{minipage}[b]{\linewidth}\raggedleft
complete\_rate
\end{minipage} & \begin{minipage}[b]{\linewidth}\raggedleft
min
\end{minipage} & \begin{minipage}[b]{\linewidth}\raggedleft
max
\end{minipage} & \begin{minipage}[b]{\linewidth}\raggedleft
empty
\end{minipage} & \begin{minipage}[b]{\linewidth}\raggedleft
n\_unique
\end{minipage} & \begin{minipage}[b]{\linewidth}\raggedleft
whitespace
\end{minipage} \\
\midrule\noalign{}
\endhead
\bottomrule\noalign{}
\endlastfoot
date & 0 & 1 & 10 & 10 & 0 & 391 & 0 \\
time & 0 & 1 & 8 & 8 & 0 & 24 & 0 \\
\end{longtable}
}

\textbf{Variable type: numeric}

{\def\LTcaptype{none} % do not increment counter
\begin{longtable}[]{@{}
  >{\raggedright\arraybackslash}p{(\linewidth - 20\tabcolsep) * \real{0.1429}}
  >{\raggedleft\arraybackslash}p{(\linewidth - 20\tabcolsep) * \real{0.1020}}
  >{\raggedleft\arraybackslash}p{(\linewidth - 20\tabcolsep) * \real{0.1429}}
  >{\raggedleft\arraybackslash}p{(\linewidth - 20\tabcolsep) * \real{0.0816}}
  >{\raggedleft\arraybackslash}p{(\linewidth - 20\tabcolsep) * \real{0.0714}}
  >{\raggedleft\arraybackslash}p{(\linewidth - 20\tabcolsep) * \real{0.0816}}
  >{\raggedleft\arraybackslash}p{(\linewidth - 20\tabcolsep) * \real{0.0816}}
  >{\raggedleft\arraybackslash}p{(\linewidth - 20\tabcolsep) * \real{0.0714}}
  >{\raggedleft\arraybackslash}p{(\linewidth - 20\tabcolsep) * \real{0.0816}}
  >{\raggedleft\arraybackslash}p{(\linewidth - 20\tabcolsep) * \real{0.0816}}
  >{\raggedright\arraybackslash}p{(\linewidth - 20\tabcolsep) * \real{0.0612}}@{}}
\toprule\noalign{}
\begin{minipage}[b]{\linewidth}\raggedright
skim\_variable
\end{minipage} & \begin{minipage}[b]{\linewidth}\raggedleft
n\_missing
\end{minipage} & \begin{minipage}[b]{\linewidth}\raggedleft
complete\_rate
\end{minipage} & \begin{minipage}[b]{\linewidth}\raggedleft
mean
\end{minipage} & \begin{minipage}[b]{\linewidth}\raggedleft
sd
\end{minipage} & \begin{minipage}[b]{\linewidth}\raggedleft
p0
\end{minipage} & \begin{minipage}[b]{\linewidth}\raggedleft
p25
\end{minipage} & \begin{minipage}[b]{\linewidth}\raggedleft
p50
\end{minipage} & \begin{minipage}[b]{\linewidth}\raggedleft
p75
\end{minipage} & \begin{minipage}[b]{\linewidth}\raggedleft
p100
\end{minipage} & \begin{minipage}[b]{\linewidth}\raggedright
hist
\end{minipage} \\
\midrule\noalign{}
\endhead
\bottomrule\noalign{}
\endlastfoot
co\_gt & 1592 & 0.83 & -0.22 & 21.80 & -200.00 & 1.00 & 1.8 & 2.90 &
11.90 & ▁▁▁▁▇ \\
pt08\_s1\_co & 366 & 0.96 & 1099.83 & 217.08 & 647.00 & 937.00 & 1063.0
& 1231.00 & 2040.00 & ▃▇▃▁▁ \\
nmhc\_gt & 8443 & 0.10 & 218.81 & 204.46 & 7.00 & 67.00 & 150.0 & 297.00
& 1189.00 & ▇▂▁▁▁ \\
c6h6\_gt & 0 & 1.00 & 1.87 & 41.38 & -200.00 & 4.00 & 7.9 & 13.60 &
63.70 & ▁▁▁▇▅ \\
pt08\_s2\_nmhc & 366 & 0.96 & 939.15 & 266.83 & 383.00 & 734.50 & 909.0
& 1116.00 & 2214.00 & ▅▇▃▁▁ \\
n\_ox\_gt & 1639 & 0.82 & 246.90 & 212.98 & 2.00 & 98.00 & 180.0 &
326.00 & 1479.00 & ▇▂▁▁▁ \\
pt08\_s3\_n\_ox & 366 & 0.96 & 835.49 & 256.82 & 322.00 & 658.00 & 806.0
& 969.50 & 2683.00 & ▇▇▁▁▁ \\
no2\_gt & 1642 & 0.82 & 113.09 & 48.37 & 2.00 & 78.00 & 109.0 & 142.00 &
340.00 & ▃▇▃▁▁ \\
pt08\_s4\_no2 & 366 & 0.96 & 1456.26 & 346.21 & 551.00 & 1227.00 &
1463.0 & 1674.00 & 2775.00 & ▂▇▇▂▁ \\
pt08\_s5\_o3 & 366 & 0.96 & 1022.91 & 398.48 & 221.00 & 731.50 & 963.0 &
1273.50 & 2523.00 & ▃▇▅▂▁ \\
t & 366 & 0.96 & 18.32 & 8.83 & -1.90 & 11.80 & 17.8 & 24.40 & 44.60 &
▂▇▇▃▁ \\
rh & 366 & 0.96 & 49.23 & 17.32 & 9.20 & 35.80 & 49.6 & 62.50 & 88.70 &
▂▆▇▆▂ \\
ah & 366 & 0.96 & 1.03 & 0.40 & 0.18 & 0.74 & 1.0 & 1.31 & 2.23 &
▃▇▇▃▁ \\
\end{longtable}
}

\textbf{Variable type: POSIXct}

{\def\LTcaptype{none} % do not increment counter
\begin{longtable}[]{@{}
  >{\raggedright\arraybackslash}p{(\linewidth - 12\tabcolsep) * \real{0.1308}}
  >{\raggedleft\arraybackslash}p{(\linewidth - 12\tabcolsep) * \real{0.0935}}
  >{\raggedleft\arraybackslash}p{(\linewidth - 12\tabcolsep) * \real{0.1308}}
  >{\raggedright\arraybackslash}p{(\linewidth - 12\tabcolsep) * \real{0.1869}}
  >{\raggedright\arraybackslash}p{(\linewidth - 12\tabcolsep) * \real{0.1869}}
  >{\raggedright\arraybackslash}p{(\linewidth - 12\tabcolsep) * \real{0.1869}}
  >{\raggedleft\arraybackslash}p{(\linewidth - 12\tabcolsep) * \real{0.0841}}@{}}
\toprule\noalign{}
\begin{minipage}[b]{\linewidth}\raggedright
skim\_variable
\end{minipage} & \begin{minipage}[b]{\linewidth}\raggedleft
n\_missing
\end{minipage} & \begin{minipage}[b]{\linewidth}\raggedleft
complete\_rate
\end{minipage} & \begin{minipage}[b]{\linewidth}\raggedright
min
\end{minipage} & \begin{minipage}[b]{\linewidth}\raggedright
max
\end{minipage} & \begin{minipage}[b]{\linewidth}\raggedright
median
\end{minipage} & \begin{minipage}[b]{\linewidth}\raggedleft
n\_unique
\end{minipage} \\
\midrule\noalign{}
\endhead
\bottomrule\noalign{}
\endlastfoot
datetime & 2 & 1 & 2004-03-10 18:00:00 & 2005-04-04 14:00:00 &
2004-09-21 16:00:00 & 9355 \\
\end{longtable}
}

\begin{Shaded}
\begin{Highlighting}[]
\FunctionTok{vis\_dat}\NormalTok{(df)}
\end{Highlighting}
\end{Shaded}

\pandocbounded{\includegraphics[keepaspectratio]{lab2_files/figure-latex/unnamed-chunk-8-1.pdf}}

\begin{Shaded}
\begin{Highlighting}[]
\FunctionTok{vis\_miss}\NormalTok{(df)}
\end{Highlighting}
\end{Shaded}

\pandocbounded{\includegraphics[keepaspectratio]{lab2_files/figure-latex/unnamed-chunk-8-2.pdf}}

\begin{Shaded}
\begin{Highlighting}[]
\NormalTok{num\_df }\OtherTok{\textless{}{-}}\NormalTok{ df }\SpecialCharTok{\%\textgreater{}\%} \FunctionTok{select}\NormalTok{(}\FunctionTok{where}\NormalTok{(is.numeric))}
\FunctionTok{vis\_expect}\NormalTok{(num\_df, }\SpecialCharTok{\textasciitilde{}}\NormalTok{ .x }\SpecialCharTok{\textless{}} \DecValTok{0}\NormalTok{)}
\end{Highlighting}
\end{Shaded}

\pandocbounded{\includegraphics[keepaspectratio]{lab2_files/figure-latex/unnamed-chunk-9-1.pdf}}

Комментарий: \texttt{vis\_expect()} подсвечивает отрицательные значения,
которые для концентраций и метеопараметров недопустимы

\begin{Shaded}
\begin{Highlighting}[]
\FunctionTok{create\_report}\NormalTok{(df, }\AttributeTok{output\_file =} \StringTok{"profile.html"}\NormalTok{)}
\end{Highlighting}
\end{Shaded}

\begin{verbatim}
##   |                                             |                                     |   0%  |                                             |.                                    |   2%                                   |                                             |..                                   |   4% [global_options]                  |                                             |..                                   |   7%                                   |                                             |...                                  |   9% [introduce]                       |                                             |....                                 |  11%                                   |                                             |.....                                |  13% [plot_intro]
\end{verbatim}

\begin{verbatim}
##   |                                             |......                               |  15%                                   |                                             |......                               |  17% [data_structure]                  |                                             |.......                              |  20%                                   |                                             |........                             |  22% [missing_profile]
\end{verbatim}

\begin{verbatim}
##   |                                             |.........                            |  24%                                   |                                             |..........                           |  26% [univariate_distribution_header]  |                                             |..........                           |  28%                                   |                                             |...........                          |  30% [plot_histogram]
\end{verbatim}

\begin{verbatim}
##   |                                             |............                         |  33%                                   |                                             |.............                        |  35% [plot_response_histogram]         |                                             |..............                       |  37%                                   |                                             |..............                       |  39% [plot_density]                    |                                             |...............                      |  41%                                   |                                             |................                     |  43% [plot_response_density]           |                                             |.................                    |  46%                                   |                                             |..................                   |  48% [plot_frequency_bar]
\end{verbatim}

\begin{verbatim}
##   |                                             |..................                   |  50%                                   |                                             |...................                  |  52% [plot_response_bar]               |                                             |....................                 |  54%                                   |                                             |.....................                |  57% [plot_with_bar]                   |                                             |......................               |  59%                                   |                                             |.......................              |  61% [plot_normal_qq]
\end{verbatim}

\begin{verbatim}
##   |                                             |.......................              |  63%                                   |                                             |........................             |  65% [plot_response_qq]                |                                             |.........................            |  67%                                   |                                             |..........................           |  70% [plot_by_qq]                      |                                             |...........................          |  72%                                   |                                             |...........................          |  74% [correlation_analysis]
\end{verbatim}

\begin{verbatim}
##   |                                             |............................         |  76%                                   |                                             |.............................        |  78% [principal_component_analysis]
\end{verbatim}

\begin{verbatim}
##   |                                             |..............................       |  80%                                   |                                             |...............................      |  83% [bivariate_distribution_header]   |                                             |...............................      |  85%                                   |                                             |................................     |  87% [plot_response_boxplot]           |                                             |.................................    |  89%                                   |                                             |..................................   |  91% [plot_by_boxplot]                 |                                             |...................................  |  93%                                   |                                             |...................................  |  96% [plot_response_scatterplot]       |                                             |.................................... |  98%                                   |                                             |.....................................| 100% [plot_by_scatterplot]           
\end{verbatim}

\begin{verbatim}
## "C:/Program Files/RStudio/resources/app/bin/quarto/bin/tools/pandoc" +RTS -K512m -RTS "C:\Users\User\1_HomeWork\1\3\Project_3\report.knit.md" --to html4 --from markdown+autolink_bare_uris+tex_math_single_backslash --output pandoc2d742f3d703b.html --lua-filter "C:\Users\User\AppData\Local\R\win-library\4.6\rmarkdown\rmarkdown\lua\pagebreak.lua" --lua-filter "C:\Users\User\AppData\Local\R\win-library\4.6\rmarkdown\rmarkdown\lua\latex-div.lua" --lua-filter "C:\Users\User\AppData\Local\R\win-library\4.6\rmarkdown\rmarkdown\lua\table-classes.lua" --embed-resources --standalone --variable bs3=TRUE --section-divs --table-of-contents --toc-depth 6 --template "C:\Users\User\AppData\Local\R\win-library\4.6\rmarkdown\rmd\h\default.html" --syntax-highlighting none --variable highlightjs=1 --variable theme=yeti --mathjax --variable "mathjax-url=https://mathjax.rstudio.com/latest/MathJax.js?config=TeX-AMS-MML_HTMLorMML" --include-in-header "C:\Users\User\AppData\Local\Temp\RtmpAvzobS\rmarkdown-str2d7411f23cc2.html"
\end{verbatim}

\begin{Shaded}
\begin{Highlighting}[]
\FunctionTok{get\_dupes}\NormalTok{(df, datetime)}
\end{Highlighting}
\end{Shaded}

\begin{verbatim}
## # A tibble: 2 x 17
##   datetime dupe_count date       time     co_gt pt08_s1_co nmhc_gt c6h6_gt
##   <dttm>        <int> <chr>      <chr>    <dbl>      <dbl>   <dbl>   <dbl>
## 1 NA                2 28/03/2004 02:00:00   1.7       1135      95     6.3
## 2 NA                2 27/03/2005 02:00:00   1         1008      NA     3.3
## # i 9 more variables: pt08_s2_nmhc <dbl>, n_ox_gt <dbl>, pt08_s3_n_ox <dbl>,
## #   no2_gt <dbl>, pt08_s4_no2 <dbl>, pt08_s5_o3 <dbl>, t <dbl>, rh <dbl>,
## #   ah <dbl>
\end{verbatim}

\begin{Shaded}
\begin{Highlighting}[]
\FunctionTok{count}\NormalTok{(df, }\AttributeTok{date =} \FunctionTok{as.Date}\NormalTok{(datetime), }\AttributeTok{sort =} \ConstantTok{TRUE}\NormalTok{) }\SpecialCharTok{\%\textgreater{}\%} \FunctionTok{head}\NormalTok{(}\DecValTok{10}\NormalTok{)}
\end{Highlighting}
\end{Shaded}

\begin{verbatim}
## # A tibble: 10 x 2
##    date           n
##    <date>     <int>
##  1 2004-03-11    24
##  2 2004-03-12    24
##  3 2004-03-13    24
##  4 2004-03-14    24
##  5 2004-03-15    24
##  6 2004-03-16    24
##  7 2004-03-17    24
##  8 2004-03-18    24
##  9 2004-03-19    24
## 10 2004-03-20    24
\end{verbatim}

Пропуск: \texttt{get\_dupes()} проверяет дубли времени

\textbf{Полнота по столбцам}

\begin{Shaded}
\begin{Highlighting}[]
\FunctionTok{skim}\NormalTok{(df) }\SpecialCharTok{\%\textgreater{}\%}
  \FunctionTok{select}\NormalTok{(skim\_variable, n\_missing, complete\_rate)}
\end{Highlighting}
\end{Shaded}

\begin{verbatim}
## # A tibble: 16 x 3
##    skim_variable n_missing complete_rate
##    <chr>             <int>         <dbl>
##  1 datetime              2        1.000 
##  2 date                  0        1     
##  3 time                  0        1     
##  4 co_gt              1592        0.830 
##  5 pt08_s1_co          366        0.961 
##  6 nmhc_gt            8443        0.0977
##  7 c6h6_gt               0        1     
##  8 pt08_s2_nmhc        366        0.961 
##  9 n_ox_gt            1639        0.825 
## 10 pt08_s3_n_ox        366        0.961 
## 11 no2_gt             1642        0.825 
## 12 pt08_s4_no2         366        0.961 
## 13 pt08_s5_o3          366        0.961 
## 14 t                   366        0.961 
## 15 rh                  366        0.961 
## 16 ah                  366        0.961
\end{verbatim}

Комментарий: столбцы с \texttt{complete\_rate} заметно ниже 1 ---
кандидаты на пропуски

\section{7. Чек-лист и план
очистки}\label{ux447ux435ux43a-ux43bux438ux441ux442-ux438-ux43fux43bux430ux43d-ux43eux447ux438ux441ux442ux43aux438}

{\def\LTcaptype{none} % do not increment counter
\begin{longtable}[]{@{}
  >{\raggedright\arraybackslash}p{(\linewidth - 10\tabcolsep) * \real{0.1667}}
  >{\raggedright\arraybackslash}p{(\linewidth - 10\tabcolsep) * \real{0.1667}}
  >{\raggedright\arraybackslash}p{(\linewidth - 10\tabcolsep) * \real{0.1667}}
  >{\raggedright\arraybackslash}p{(\linewidth - 10\tabcolsep) * \real{0.1667}}
  >{\raggedright\arraybackslash}p{(\linewidth - 10\tabcolsep) * \real{0.1667}}
  >{\raggedright\arraybackslash}p{(\linewidth - 10\tabcolsep) * \real{0.1667}}@{}}
\toprule\noalign{}
\begin{minipage}[b]{\linewidth}\raggedright
№
\end{minipage} & \begin{minipage}[b]{\linewidth}\raggedright
Где (столбец, строки)
\end{minipage} & \begin{minipage}[b]{\linewidth}\raggedright
Проблема
\end{minipage} & \begin{minipage}[b]{\linewidth}\raggedright
Чем подтверждено
\end{minipage} & \begin{minipage}[b]{\linewidth}\raggedright
Измерение качества
\end{minipage} & \begin{minipage}[b]{\linewidth}\raggedright
Предлагаемое решение
\end{minipage} \\
\midrule\noalign{}
\endhead
\bottomrule\noalign{}
\endlastfoot
1 & числовые датчики & скрытый код пропуска \texttt{-200} &
\texttt{read\_lines()}, значения \texttt{-200} & полнота & считать
\texttt{-200} как \texttt{NA} при импорте \\
2 & весь файл & десятичная запятая & \texttt{read\_lines()},
\texttt{read\_csv2()} & структура/типы & использовать
\texttt{read\_csv2()} или \texttt{locale(decimal\_mark\ =\ ",")} \\
3 & \texttt{Date}, \texttt{Time} & дата и время хранятся текстом &
\texttt{glimpse()} & типы/согласованность & собрать \texttt{POSIXct}
через \texttt{parse\_date\_time()} \\
4 & конец файла & пустые строки и столбцы &
\texttt{read\_lines(skip\ =\ 9355)}, \texttt{problems()} & структура &
\texttt{remove\_empty()} \\
5 & имена столбцов & несинтаксические имена \texttt{CO(GT)} &
\texttt{glimpse()} & структура & \texttt{clean\_names()} \\
6 & числовые столбцы & отрицательные значения там, где они невозможны &
\texttt{vis\_expect(\textasciitilde{}\ .x\ \textless{}\ 0)} & валидность
& заменить на \texttt{NA}/проверить источник \\
7 & \texttt{co\_gt}, \texttt{nmhc\_gt}, \texttt{nox\_gt} и др. & выбросы
и широкие диапазоны & \texttt{skim()} min/max/quantiles &
точность/валидность & проверить по доменным границам \\
8 & \texttt{datetime} & не указан часовой пояс &
\texttt{parse\_date\_time(tz=...)} & согласованность & явно задать
\texttt{tz\ =\ "Europe/Rome"} \\
9 & \texttt{datetime} & возможны дубли меток времени &
\texttt{get\_dupes(df,\ datetime)} & уникальность & решить, агрегировать
или оставить \\
10 & \texttt{date}, \texttt{time} & возможны пропущенные часы &
\texttt{count(df,\ date)} & полнота/timeliness & проверить регулярность
ряда \\
\end{longtable}
}

Вывод: Наиболее серьёзны скрытые пропуски \texttt{-200} и отсутствие
единой метки времени --- они напрямую искажают статистику и не дают
строить временные ряды. Также важны отрицательные значения и выбросы
датчиков: без доменных границ их нельзя отличить от реальных измерений.
Дубли и пропущенные часы критичны для временного анализа.

\section{Контрольные
вопросы}\label{ux43aux43eux43dux442ux440ux43eux43bux44cux43dux44bux435-ux432ux43eux43fux440ux43eux441ux44b}

\begin{enumerate}
\def\labelenumi{\arabic{enumi}.}
\tightlist
\item
  Проверял полноту (\texttt{skim()}, \texttt{vis\_miss()}), типы и
  структуру (\texttt{glimpse()}, \texttt{vis\_dat()}), валидность
  диапазонов (\texttt{skim()}, \texttt{vis\_expect()}), уникальность
  (\texttt{get\_dupes()}), согласованность даты/времени, кодировку
  (\texttt{guess\_encoding()}), проблемы импорта (\texttt{problems()}).
\item
  \texttt{readr} угадывает тип по первым строкам; если там пропуски/коды
  вроде \texttt{-200} или пустые строки, тип может стать character. Надо
  явно задавать \texttt{col\_types} и \texttt{na}.
\item
  Скрытые коды \texttt{-200} не распознаются как \texttt{NA} и портят
  min/max/mean/квантили, графики и модели.
\item
  \texttt{get\_dupes()} только показывает дубликаты, потому что
  одинаковые строки не всегда ошибочны; решение об удалении зависит от
  контекста.
\item
  Нарушение tidy data --- это неверная форма (переменные в
  столбцах/строках, несколько значений в одной ячейке). Некорректное
  значение --- это ошибка содержания при правильной форме.
\end{enumerate}

\end{document}
